\chapter{Harvard Thinking: How to Tell a Story}

This chapter distills a panel conversation among James Wood, Sam Marks, Lauren Groff, and Nick White, moderated by Samantha Liney-Perfoss, curated in this collection by LazyingArt LLC. There is no board here, and so there is no literal classroom derivation to recover. Still, the lecture has a real spine. It begins by imposing a severe standard on writing, then moves through origins, character and plot, the writer's own obsessions and body, failure and revision, and finally the reader's test of vitality. We will follow that order closely. Whenever the discussion suggests a clean relation, we will write it schematically, but we should remember that these formulas are clarifications of the lecture, not laws written on a blackboard.

\section{Opening Criterion: Writing That Costs Something}

The lecture opens unusually hard. It does not begin with character, plot, or craft terminology. It begins by asking what sort of writing can affect another person at all. The answer is that emotionally effective writing must itself be written emotionally. That is not a piece of encouragement. It is a criterion.

In compact form, the opening claim is
\begin{equation}
\text{emotional effect on reader} \Rightarrow \text{emotional investment by writer}.
\end{equation}
The lecture immediately strengthens this. What matters is not simply that time was spent at the desk. Writing that works costs more than time. It asks for strain, exposure, and risk.
\begin{equation}
\text{working writing} \Rightarrow C > T,
\end{equation}
where \(T\) is mere time spent and \(C\) is the wider cost of the work. We should not read this as a quantitative inequality. It is a way of saying that live writing demands more than attendance.

The same point is restated in the image of ``all cylinders firing.'' When something is working, the writer feels pushed to the edge of his or her emotional and intellectual capacities. Schematically,
\begin{equation}
\text{working story} \Rightarrow E + I,
\end{equation}
with \(E\) for emotional capability engaged and \(I\) for intellectual capability engaged. Again, the notation is schematic. The lecture is telling us that good writing is not purely felt and not purely thought; it requires both, active at once.

Only after this strong opening does the host widen the frame. Storytelling is part of human life. There is the familiar image of the writer wrestling with the page. There is, certainly, craft. But the lecture does not want us to mistake craft for the whole story. The real question therefore appears at exactly the right time: what makes a story unforgettable?

\subsection{Question \& Answer}

\paragraph{Question.} What is the lecture's first test for whether writing is alive?

\paragraph{Answer.} The first test is not polish by itself. It is whether the writing has cost the writer something beyond time, and whether emotional and intellectual effort have both been brought fully into play. The lecture begins by insisting that aliveness on the page requires aliveness in the act of writing.

\section{Where a Story Begins}

Once the standard is set, the host asks the first practical question: where does a story start? The panel answers by refusing a single origin. For one writer, the seed may be carried for years. For another, the story does not really exist until writing begins. For another still, a beginning or ending line can orient the whole structure before the middle is visible.

One especially clear account describes a long-incubated idea that remains inert until it collides with some later event, reading, or experience. Only then does it acquire urgency enough to become story material. We may summarize that relation as
\begin{equation}
L + X \Rightarrow U + D + G + W,
\end{equation}
where \(L\) is a long-held idea, \(X\) a later collision, and \(U,D,G,W\) denote urgency, density, gravity, and weight. The lecture's emphasis falls on the transition from mere possession of material to pressure strong enough to demand form.

A second account is complementary rather than contradictory. The story is not simply stored in advance waiting to be transcribed. It becomes itself in the act of writing:
\begin{equation}
\text{incubation} \to \text{first writing} \to \text{unfolding}.
\end{equation}
This is one of the lecture's most important corrections to any naive model of inspiration. The story is not always fully known beforehand. Often it becomes legible only while being written.

A third account introduces orientation. Even if the middle is unclear, a first line, a beginning, or an ending can be enough to stabilize the space:
\begin{equation}
(\text{known beginning}) \lor (\text{known ending}) \Rightarrow \text{narrative orientation}.
\end{equation}
That point matters because it explains why very different writers may all feel that they know ``where to start'' even when what they mean by that phrase is quite different.

The most practical engine in this part of the lecture comes from character. One writer says that writing often begins with a particular person, and then with two immediate questions: what does that person want, and what stands in the way?
\begin{equation}
\text{story engine} = \text{character} + \text{want} + \text{obstacle}.
\end{equation}
This is a genuinely useful local formula. It already contains pressure, and pressure is what gives narrative motion.

\paragraph{A Worked Example.} We can combine several of the panel's origin accounts without pretending that they form a universal algorithm:
\begin{align}
L &\text{ may remain inert for a long time},\\
L + X &\Rightarrow U + D + G + W,\\
U + D + G + W &\Rightarrow \text{first writing},\\
\text{first writing} &\Rightarrow \text{unfolding}.
\end{align}
This is not a law of storytelling. It is a careful reconstruction of the lecture's shared pattern: incubation, activation, first writing, and emergence.

The section closes with an important negative lesson. The writers warn against the fantasy that one can rescue a weak draft by arbitrary twist. The moment one thinks, ``and then this crazy thing happens,'' the result is often disaster. The origin of a story must feel earned, not merely decorative.

\subsection{Question \& Answer}

\paragraph{Question.} Where does a story actually begin?

\paragraph{Answer.} The lecture's answer is plural. A story may begin in long incubation, in the collision of an old idea with a new experience, in a first line, in an ending, or in a character under pressure. But in several accounts the decisive moment comes when those materials pass into writing and begin to unfold there.

\section{Character, Action, and the Plot Spectrum}

The lecture now narrows its focus. Once character has been introduced as an engine, the next question is what relation character bears to action. The discussion first stays close to experience. Characters may surprise the writer, reveal themselves gradually, or seem clearer in some scenes than in others. In dramatic writing they may be heard more than seen. Two characters placed together in a scene can behave like elements in an experiment: one watches what sparks.

That experimental language is important, because it leads directly into a theoretical pole. Aristotle is invoked as the representative of the classical view in which character is subordinate to action:
\begin{equation}
\text{character} \prec \text{action}.
\end{equation}
The panel does not adopt this as a full modern program. It uses it instead as one end of the field. The contemporary problem is not how to obey that principle absolutely, but how to think about a narrative space in which character and plot are both active.

That space is described as a spectrum:
\begin{equation}
\text{character-driven} \leftrightarrow \text{plot-driven}.
\end{equation}
The value of the spectrum is diagnostic. If plot dominates too aggressively, we get false propulsion, arbitrary twist, and ``plotting for the sake of plotting.'' If character dominates without story pressure, the reader may begin to feel that nothing is happening. We can therefore write the two failure modes as
\begin{equation}
\text{dishonest shoehorning} \;\;\leftrightarrow\;\; \text{character without story pressure}.
\end{equation}
The lecture's central claim is that the writer must avoid both.

The comic spaceship example makes the point sharply. Sometimes a writer introduces a wild event not because the story has earned it, but because the writer wants rescue from the difficulty of discovering what really follows. The lecture's response is subtle: sometimes the spaceship does not belong at page eight at all. Sometimes it marks the place where the story should actually begin. That is a beautiful correction, because it turns a bad escape into a question of structure.

The section also insists that character is not a fixed interior nugget. Environment matters. The same pair of people having the same breakup will not produce the same scene everywhere:
\begin{align}
(\text{same couple},\text{ beach}) &\Rightarrow \text{one field of behavior},\\
(\text{same couple},\text{ Everest}) &\Rightarrow \text{another field of behavior}.
\end{align}
Thus character is not only what a person ``is.'' It is also what that person becomes under circumstance, place, pressure, and world.

The lecture then adds a striking prompt: have your characters do the thing you are afraid to do. This is the pivot. We have moved from narrative mechanics to the writer's deeper relation to the work.

\subsection{Question \& Answer}

\paragraph{Question.} How should we think about character and plot without flattening one into the other?

\paragraph{Answer.} The lecture treats them as mutually necessary forms of pressure. Character gives us desire, obstacle, and response; plot gives that pressure direction and consequence. The failure comes when plot becomes arbitrary or when character loses narrative force. Good storytelling keeps both alive at once.

\section{Soul, Obsession, Place, and the Body}

At this point the host widens the inquiry again. If the lecture has spoken of character-driven and plot-driven writing, might there also be something like soul-of-the-writer-driven writing? This is the right moment for that question, because the previous section has already suggested that the writer's own fear and courage may be at work inside the characters.

The answers move toward obsession, place, and repeated inward questions. One writer speaks of an early novel as author-driven without being literally autobiographical. Another speaks of writing different works that are, at a deeper level, different versions of the same inward problem. Place enters with unusual force. A region such as Mississippi is not merely a setting chosen from outside; it is something lodged in the bones. One may spend years trying to leave it and then spend years trying to write back into it.

The strongest formal correction in this section is
\begin{equation}
\text{fictional distance} \not\Rightarrow \text{absence of self}.
\end{equation}
A story can be remote from a writer's lived circumstances and still be deeply personal. This matters because the lecture wants to protect fiction against two opposite misunderstandings: that the personal must be literal, and that invention must be impersonal.

The bridge into distant material is bodily rather than merely biographical. Sensory memory allows access:
\begin{equation}
\text{sensory memory} \Rightarrow \text{access to distant fictional world}.
\end{equation}
The lecture's concrete example is perfect. One need not have lived in a commune to know the feel of a handmade bowl of warm oatmeal in the hands. That is already bodily knowledge, and bodily knowledge can be transferred.

The panel's language here becomes more metaphorical: obsessions, the deepest part of the self, the writer's soul, each character as a prismatic hologram, the titration of closeness offered to the reader. We should preserve that language, but not misread it. The claim is not mystical vagueness. It is that fiction draws on recurring structures of feeling and memory which survive translation into new forms, new periods, and new settings.

There is also an important practical consequence. If the writer returns again and again to certain places, fears, or questions, this need not be failure of invention. It may be the sign that the work has found its real field of pressure.

\subsection{Question \& Answer}

\paragraph{Question.} How can fiction be deeply personal without being literally autobiographical?

\paragraph{Answer.} The lecture's answer is that the self enters fiction through obsession, place, bodily memory, fear, desire, and recurring inward questions. These do not require literal self-report. A writer can invent outward circumstance and still be working from the deepest available materials.

\section{Failure, Drafting, Revision, and the Reader}

The lecture now reaches its most detailed section on process. Failure is not treated as an embarrassment to be hidden after the fact. It is treated as a working medium. One writer openly declares love of failure, not because failure is pleasant in itself, but because it permits play, restarting, and genuine experiment without premature cleanliness.

The lecture's process relation can therefore be written as
\begin{equation}
\text{draft} \to \text{failure} \to \text{restart} \to \text{story teaches itself}.
\end{equation}
The logic is precise. Drafting produces imperfect material. Imperfection is not simply endured; it is used. Restarting is not mere repetition, but another attempt under clearer constraints. After enough passes, the story begins to show its own demands more distinctly.

The panel then compares working methods. One writer never rereads notebooks and remains purely analog almost until the end. Another moves between legal pads, typed pages, printouts, and corrections. Another emphasizes carrying the story through daily life, so that writing is not only what happens at the desk but what persists while walking, worrying, and trying to untie the knots one has written oneself into. These methods differ, but the lecture does not make them compete as dogmas. Their shared point is that process must be deliberate and that different processes produce different kinds of contact with the material.

A particularly strong insight arrives when revision is linked to the reader. Drafting in heat can make bad work seem brilliant because it is still charged with the writer's own immediacy. Revision introduces another mind. Hence the beginning of a story is not merely the first temporal portion of the work:
\begin{equation}
\text{beginning} \Rightarrow \text{teaches reader how to read the story}.
\end{equation}
This is why beginnings are slow. They are not only openings; they are instructions.

\paragraph{Derivation of the Revision Principle.} The lecture's logic can be stated step by step.
\begin{enumerate}
\item Early drafting is largely private; the writer is, for the moment, the only reader.
\item Revision introduces another reader, even if only an imagined one.
\item The opening paragraphs are the first place where that reader learns the text's rhythm, scale, and mode of attention.
\item Therefore revision bears unusually hard on beginnings.
\item Once the beginning has taught the reader correctly, the rest of the motion can proceed more freely.
\end{enumerate}

The lecture also insists on temporal sobriety. In the heat of drafting, we may sincerely believe we have written something brilliant. The next day, with distance restored, the same passage may look very different. That is not a betrayal of inspiration. It is the necessary arrival of judgment.

\subsection{Question \& Answer}

\paragraph{Question.} What does failure do in the writing process?

\paragraph{Answer.} Failure reveals the limits of the current attempt. Once it is accepted as part of the work rather than as humiliation outside the work, it becomes productive. It drives restarting, clarifies demand, and gradually allows the story to teach the writer what shape it wants.

\section{Vitality, Reading, Editing, and Final Advice}

The final movement asks what makes a story good. The first answer is negative in a productive way. Goodness does not lie first in whether the story is realist or experimental, continuous or fragmented, conventional or strange. It lies first in vitality:
\begin{equation}
\text{good story} \Rightarrow \text{vitality} = \text{non-deadness} = \text{being drawn in}.
\end{equation}
That is the lecture's late governing criterion, and it answers to the opening criterion of cost. The work that has cost the writer something should produce life on the page.

The lecture now turns to reading. Reading itself is divided into two stages:
\begin{equation}
\text{honeymoon read} \to \text{seven-year itch}.
\end{equation}
The first reading suspends judgment and lets the piece act on us. The second rereads analytically, asks why a passage mattered, and starts assembling an argument about the work. This is offered both as a critic's method and as a workshop method.

The writer's own analogue of vitality is a bodily criterion. When the work is approaching what it wants to be, one can feel it:
\begin{equation}
\text{comes close} \Rightarrow \text{you can feel it}.
\end{equation}
The lecture calls this a vibration. That is metaphorical language, but it is not idle language. It means that some kinds of formal success are first known by contact before they are known by explicit theory.

At exactly this point, the lecture returns to its opening lines. We are told again that emotionally effective writing must be done emotionally, that it costs more than time, and that it presses the writer to the edge of available capability. This return matters. It closes the circle. Vitality on the page is not detached from the cost of making it.

The next practical test is reading aloud. Reading aloud estranges the text from us and turns us into colder judges of it:
\begin{equation}
\text{reading aloud} \Rightarrow \text{estrangement} \Rightarrow \text{bullshit detector}.
\end{equation}
A sentence that looked acceptable in silent rereading may fail instantly in public speech. The lecture's point is that such failure should not be explained away. It is evidence.

This leads naturally into ``murder your darlings.'' But again the lecture resists a slogan where a judgment is required:
\begin{equation}
\text{edit well} \neq \text{cut mechanically}.
\end{equation}
Some material truly is indulgent and must go. Other material is difficult because it has not yet been written well enough. That is why the best editorial mark in the lecture is not ``remove this,'' but ``do better here.'' Good editing identifies places where life is weak and asks for a stronger sentence, not merely a shorter manuscript.

The closing advice translates everything into habit. Stay in contact with the work:
\begin{equation}
\text{touch the work daily} \Rightarrow \text{continued contact with the story}.
\end{equation}
Writing is first of all an activity, not merely a thing one later possesses:
\begin{equation}
\text{writing} = \text{process} \quad \text{before it is} \quad \text{product}.
\end{equation}
And finally, one should not become hidebound by imported rules. The lecture ends by saying, in effect, that the only rules which finally matter are those generated by the story itself:
\begin{equation}
\text{rules} = \text{rules inherent to the story at hand}.
\end{equation}
This is not a permission slip for laziness. It is a demand for greater attention. The story must be listened to closely enough for its own rules to become audible.

\subsection{Question \& Answer}

\paragraph{Question.} What makes a story work, and how do we know when revision is helping rather than deadening it?

\paragraph{Answer.} A story works when it has vitality, when it feels alive and draws us in. Revision helps when it increases that life, whether by removing indulgence, exposing false notes, or strengthening passages that have not yet earned themselves. Revision fails only when it becomes mechanical rule-following rather than renewed attention to where the work is most alive.

\section{Summary}

The lecture begins with cost and ends with vitality. Between those two points it unfolds a coherent account of storytelling. Stories may incubate for years, may come alive through collision, may begin in character or in a first line, and may only truly emerge in the act of writing. Character and plot must be balanced; the self enters fiction through obsession, place, and bodily memory; failure is not external to the work but part of its medium; revision introduces the reader; and judgment arrives through rereading, reading aloud, and editorial distance.

If we compress the whole lecture to one final sequence, it is this:
\begin{equation}
\text{pressure} \to \text{writing} \to \text{failure} \to \text{revision} \to \text{vitality}.
\end{equation}
That chain is not everything, but it is the line of force running through the discussion. A story begins in pressure, takes shape in language, survives failure, learns from revision, and is finally tested by whether it lives on the page.